\documentclass{article}
\usepackage[utf8]{inputenc}
\usepackage[margin=1in,left=1.5in,includefoot]{geometry}
\usepackage{booktabs}
\usepackage{float}
% Header & Footer Stuff

\usepackage{fancyhdr}
\usepackage{wasysym}
\usepackage{amsmath}
\usepackage{tikz}
\usepackage{tabularx}
\usepackage{amssymb}
\usepackage{graphicx}
\pagestyle{fancy}
\lhead{MATH 220/199}
\rhead{May 7, 2026}
\fancyfoot{}
\lfoot{David Gallardo}
\fancyfoot[R]{}
%027601193
% The Main Document
\begin{document}

\begin{center}
 \LARGE\bfseries MATH 220/199 Practice Final Exam
 \line(1,0){430}
\end{center}
Name:
\section*{Instructions}
\begin{itemize}
    \item Please write your name.
    \item Show as much work as you can for every question.
    \item Attempt every single question.
    \item Attempt every single question.
    \item I cannot emphasize this enough - Attempt every single question.
    \item The use of electronics - including a calculator - is prohibited. Leave your answers in exact form if you must.
    \item Any form of cheating is strictly prohibited.
    \item GOOD LUCK!!!
\end{itemize}
\textbf{REMARK:} This practice exam is not necessarily comprehensive, but every topic in this practice exam is fair game for the actual exam.
\newpage
\section*{Chapter 1: Functions and Models}
\begin{enumerate}
    \item Let $f(x) = x^2-\sqrt{2x}, g(x) = \cos(2x)$.
    \begin{enumerate}
        \item Compute $f(g(x))$
        \item Compute $g(f(x))$
    \end{enumerate}
    \vfill
    \item Given that $\sin(\theta) = \frac{3}{\sqrt{10}}$ Compute the following quantities:
    \begin{enumerate}
        \item $\cos(\theta)$
        \item $\cot(\theta)$
        \item $\sin(2\theta)$
    \end{enumerate}
    \vfill
\end{enumerate}
\newpage
\section*{Chapter 1: Functions and Models (Contd)}
\begin{enumerate}
    \item In the space below, make a sketch of $f(x) = 2\ln(x-2)$. Label any asymptotes, any x-intercepts, and one non-x-intercept point on the graph.
    \vfill
    \item Let $g(x) = 3e^{5x-1}+2$. Compute $g^{-1}(x)$. State the domain and range of $g^{-1}(x)$.
    \vfill
\end{enumerate}
\newpage
\section*{Chapter 2: Limits and Derivatives}
\begin{enumerate}
    \item Compute $\lim\limits_{x \to 9} \frac{x-9}{\sqrt{x}-3}$
    \vfill
    \item Compute $\lim\limits_{x \to -\infty} \frac{\sqrt{2x^2+1}}{4x}$
    \vfill
\end{enumerate}
\newpage
\section*{Chapter 2: Limits and Derivatives (Contd)}
\begin{enumerate}
    \item Given $f(x) = 3x^2+2x$ Use the limit definition of the derivative to compute $f'(x)$.
    \vfill
    \item Consider $f(x)$ which is sketched below:
    \newline
    \newline
    \includegraphics[]{pfe1.png}
    \newline
    \newline
    Compute the following or write DNE if the value does not exist: 
    \begin{enumerate}
        \item $f(1)$
        \item $\lim\limits_{x \to 1} f(x)$
        \item $\lim\limits_{x \to 3^-} f(x)$
        \item $f(2)$
    \end{enumerate}
    \vfill
\end{enumerate}
\newpage
\section*{Chapter 3: Differentiation Rules}
For each of the following functions, compute the derivative:
\begin{enumerate}
    \item $f(x) = \sin^{-1}(e^{5x^2+5x})$
    \vfill
    \item $g(x) = \frac{\ln(2x)}{5^x}$
    \vfill
    \item $h(x) = (\sqrt{x})^3\cdot \tan(x)$
    \vfill
\end{enumerate}
\newpage
\section*{Chapter 3: Differentiation Rules (Contd)}
\begin{enumerate}
    \item A cylinder with constant height 10 ft. has a radius which is increasing at a rate of 2 ft/min. When the volume of the cylinder is $40 \pi \text{ ft}^3$, how fast is the surface area of the cylinder increasing?
    \vfill
    \item Using an appropriate $f(x)$ and the Method of Linearization, find an approximation of $(2.01)^3$.
    \vfill
\end{enumerate}
\newpage
\section*{Chapter 4: Applications of Differentiation}
\begin{enumerate}
    \item Let $f(x) = \frac{1}{x}+4x^2$
    \begin{enumerate}
        \item Compute all the asymptotes of $f(x)$.
        \item Compute the x-intercepts and y-intercepts of $f(x)$.
        \item Compute the critical points of $f(x)$.
        \item Use (a), (b), and (c) to make a sketch of $f(x)$ in the space below.
    \end{enumerate}
\end{enumerate}
\newpage
\section*{Chapter 4: Applications of Differentiation (Contd)}
\begin{enumerate}
    \item Compute the following limit: $\lim\limits_{x \to \infty} \frac{3x^2+7x}{e^{2x}}$
    \vfill
    \item Prove the following statement: In the interval $(-\frac{\pi}{4}, \frac{3\pi}{4}),$ the function $g(x) = \sin(x)+\cos(x)$ has a critical point.
    \vfill
    \item A fence for feral koalas is being built, and it must have an enclosed area of $400 \text{ ft}^2$. The north and south sides of the fence are being built with material which costs 20 dollars per foot, and the east and west sides are being built with material which costs 10 dollars per foot. Find the dimensions of the fence which minimizes the total cost.
    \vfill
\end{enumerate}
\newpage
\section*{Chapter 5: Integrals}
\begin{enumerate}
    \item Given: $f(x)$ is an odd function and $\int\limits_{-2}^7 f(x) \text{ d}x = 32$, compute $\int\limits_{7}^2 f(x) \text{ d}x$.
    \vfill
    \item Compute $\frac{d}{dx} \int\limits_{31}^{x^2-x} \ln(3t^2+1) \text{ d}t$
    \vfill
    \item Compute $\int\limits_{1}^4 4x\sqrt{x^2+1} \text{ d}x$
    \vfill
\end{enumerate}
\newpage
\section*{Chapter 5: Integrals (Contd)}
\begin{enumerate}
    \item An ultra-marathoner runs at a rate given by $v(t) = 3t^2-t+2$ (in miles per hour). At hour 2, the runner is already 10 miles into the course. How far will the runner be at hour 9?
    \vfill
    \item Use the limit definition of the integral to compute $\int\limits_{-1}^3 3x-4 \text{ d}x$. (HINT: You can use usual integration to check your answer).
    \vfill
\end{enumerate}
\newpage
\section*{Chapter 6: Applications of Integration}
\begin{enumerate}
    \item Let $R$ be the region bounded between $f(x) = x^2+x$ and $g(x) = 2x+2$. Compute the area of $R$.
    \vfill
    \item Let $K$ be the region enclosed by $f(x) = \sqrt{x-2}, x=5, y=0$.
    \begin{enumerate}
        \item Setup but do not evaluate an integral to find the volume of the figure obtained by rotating $K$ around the y-axis.
        \item Setup but do not evaluate an integral to find the volume of the figure obtained by rotating $K$ around the x-axis.
    \end{enumerate}
    \vfill
\end{enumerate}
\newpage
\section*{Challenge Problems}
\begin{enumerate}
    \item Determine the set of all $n \in \mathbb{Z}$ for which $\int\limits_{0}^\pi \sin(nx) \text{ d}x = 0$.
    \vfill
    \item Let $f(x) = \sin(-x)+e^{x\sqrt[3]{3}}+\ln(x\sqrt[7]{7})$. Let $f^{(n)}$ denote the $n^{th}$ derivative of $f(x)$. What is the smallest $n$ for which $f^{(n)}$ has integer coefficients for every term?
    \vfill
\end{enumerate}
\newpage
\section*{Page for Scratch Work + Remarks}
\subsection*{Remarks}
The following topics were not covered on this practice exam but are fair game for the final. You may want to revisit worksheets or quizzes which covered the following topics/techniques:
\begin{enumerate}
    \item Chapter 1:
    \begin{enumerate}
        \item Trigonometry
    \end{enumerate}
    \item Chapter 2:
    \begin{enumerate}
        \item Continuity
    \end{enumerate}
    \item Chapter 3:
    \begin{enumerate}
        \item Expect harder material on the final regarding related rates, trigonometric derivatives, logarithmic derivatives, and linearization.
    \end{enumerate}
    \item Chapter 4:
    \begin{enumerate}
        \item Curve Sketching
        \item Anti-derivatives
    \end{enumerate}
    \item Chapter 5:
    \begin{enumerate}
        \item Make sure to review Exam 3 when it is published to see what you missed.
    \end{enumerate}
    \item Chapter 6:
    \begin{enumerate}
        \item Make sure to review Exam 3 when it is published to see what you missed.
        \item More comprehensive practice which covers more cases can be found here:
        \newline
        https://tutorial.math.lamar.edu/Problems/CalcI/VolumeWithRings.aspx
    \end{enumerate}
\end{enumerate}
\section*{Space for Scratch Work}
\end{document}
